\documentclass[a4paper,11pt]{article}
\usepackage[margin=1in]{geometry}
\usepackage{amsmath,amssymb,mathtools}
\usepackage{booktabs}
\usepackage{siunitx}

\title{Unified Atomic Synthons and BDPCS3}
\author{Merlino 3.0}
\date{\today}

\begin{document}
\maketitle

\section{Goal}
Atomic synthons are continuous, geometry-derived features used to:
\begin{enumerate}
\item define atom types in a differentiable way;
\item explore chemical spaces through descriptor manifolds;
\item correct force-field or semi-empirical geometric parameters using local chemistry-aware features.
\end{enumerate}

A first application is BDPCS3: only bond lengths are corrected from DPCS3, since the other geometric parameters are already sufficiently accurate in practice.

\section{Continuous Topology Layer}
Given Cartesian geometry $\{\mathbf{R}_i\}$ and atomic numbers $\{Z_i\}$:

\subsection{Continuous coordination number (CNA)}
For atom $i$:
\begin{equation}
\mathrm{CNA}_i = \sum_{j\neq i} g_{ij},
\qquad
g_{ij} = \frac{1}{2}\left[1+\operatorname{erf}\!\left(\alpha_{\mathrm{CNA}}\,(R^0_{ij}-R_{ij})\right)\right],
\end{equation}
with
\begin{equation}
R_{ij}=\lVert \mathbf{R}_i-\mathbf{R}_j\rVert,
\qquad
R^0_{ij}=R^{\mathrm{cov}}(Z_i)+R^{\mathrm{cov}}(Z_j).
\end{equation}
In code, $\alpha_{\mathrm{CNA}}=\texttt{CNA\_ALPHA}=8.0$.

\subsection{Coordination-dependent effective covalent radius}
Using the Pyykko table $R^{\mathrm{Pyy}}(Z,\mathrm{CN})$, define:
\begin{equation}
\tilde R^{\mathrm{cov}}_i = \mathcal{H}\!\left(Z_i,\mathrm{CNA}_i\right),
\end{equation}
where $\mathcal{H}$ is a $C^1$ Hermite interpolation between tabulated coordination radii.

\subsection{Topology bond order}
\begin{equation}
\mathrm{BO}_{ij}^{\mathrm{topo}}
=\exp\!\left(
\frac{\tilde R^{\mathrm{cov}}_i+\tilde R^{\mathrm{cov}}_j-R_{ij}}{\lambda_{\mathrm{BO}}}
\right),
\end{equation}
with $\lambda_{\mathrm{BO}}=\texttt{BO\_DECAY}=0.20$.

\section{Atomic Synthons}
Synthons are built from the same CNA/BO definitions above (uniform model).

\subsection{Electronic domains}
For main-group atoms:
\begin{align}
N_{\pi,i} &= \sum_{j\in\mathcal{N}(i)} \max(\mathrm{BO}_{ij}-1,0),\\
N_{\mathrm{res},i} &= n^{\mathrm{val}}_i-\mathrm{CNA}_i-N_{\pi,i},\\
N_{\mathrm{OS},i} &= \mathrm{mod}\!\left(\mathrm{round}(N_{\mathrm{res},i}),2\right),\\
N_{\mathrm{LP},i} &= \max\!\left(0,\frac{N_{\mathrm{res},i}-N_{\mathrm{OS},i}}{2}\right),\\
N_{\mathrm{ED},i} &= \mathrm{CNA}_i+N_{\mathrm{LP},i}+N_{\mathrm{OS},i}.
\end{align}

\subsection{Additional continuous descriptors}
For atom $i$ and bonded neighbors $\mathcal{N}(i)$:
\begin{align}
q_i &= \sum_{j\in\mathcal{N}(i)}
\frac{\chi_j-\chi_i}{(n_i+n_j)\,\max(\mathrm{BO}_{ij},\mathrm{BO}_{\min})},\\
C_i &= \frac{1}{|\mathcal{N}(i)|}\sum_{j\in\mathcal{N}(i)}
\frac{\tilde R^{\mathrm{cov}}_i+\tilde R^{\mathrm{cov}}_j}
{R_{ij}+\tilde R^{\mathrm{cov}}_i+\tilde R^{\mathrm{cov}}_j},\\
D_i &= \frac{1}{|\mathcal{N}(i)|\,\bar R_i}\sum_{j\in\mathcal{N}(i)}\left|R_{ij}-\bar R_i\right|,
\quad
\bar R_i=\frac{1}{|\mathcal{N}(i)|}\sum_{j\in\mathcal{N}(i)}R_{ij},\\
\mathrm{EAN}_i &= \frac{\sqrt{\sum_{k=1}^{4}\left[\operatorname{erf}(x_{k,i}/\sigma_{\mathrm{EAN}})\right]^2}}
{1+\sqrt{\sum_{k=1}^{4}\left[\operatorname{erf}(x_{k,i}/\sigma_{\mathrm{EAN}})\right]^2}},
\end{align}
with $x_{k,i}\in\{q_i,C_i,D_i,S_i\}$ (charge, covalency, delocalization, strain), and
\begin{equation}
Z^{\mathrm{eff}}_i = Z_i-\frac{1}{2}+\mathrm{EAN}_i.
\end{equation}

A compact atom signature can be defined as:
\begin{equation}
\Sigma_i=\left(Z_i,\ \mathrm{aromatic}_i,\ \mathrm{round}(N_{\mathrm{ED},i}),\ \mathrm{round}(\mathrm{CNA}_i)\right).
\end{equation}

\section{BDPCS3 as a Synthon-Driven Correction}
\subsection{Rationale}
BDPCS3 corrects DPCS3 bond lengths using local descriptors. In the unified implementation,
the bond-order term used by BDPCS3 is taken from topology:
\begin{equation}
\mathrm{BO}_{ij}^{\mathrm{BDPCS3}} \leftarrow \mathrm{BO}_{ij}^{\mathrm{topo}}.
\end{equation}
This enforces consistency between topology/synthons and geometry correction.

\subsection{Updated BDPCS3 equation}
For bond $(i,j)$ with DPCS3 length $r_{ij}$:
\begin{equation}
r_{ij}^{\mathrm{BDPCS3}} = r_{ij}^{\mathrm{DPCS3}}+\Delta r_{ij},
\end{equation}
\begin{equation}
\Delta r_{ij}
=A_{ij}(r_{ij})
\left[1-\chi_{\mathrm{CS}}(i)\chi_{\mathrm{CS}}(j)
\exp\!\left(-\left(\frac{r_{ij}-r^{\mathrm{cov}}_{ij}+\Delta r}{\sigma}\right)^2\right)\right],
\end{equation}
with
\begin{equation}
A_{ij}(r)
=-K\,r^{\mathrm{cov}}_{ij}
\left(1+0.25\,\chi_{\mathrm{CH},ij}\right)
\left[\max(n_i,n_j)^2-1\right]
f_{\mathrm{coord}}(r),
\end{equation}
\begin{equation}
f_{\mathrm{coord}}(r)=\frac{1}{2}\left[1-\operatorname{erf}\!\left(\frac{r-(r^{\mathrm{cov}}_{ij}+0.35)}{\sigma}\right)\right].
\end{equation}

Element selectors:
\begin{align}
\chi_{\mathrm{CS}}(x)&=\delta_{Z_x,6}+0.5\,\delta_{Z_x,16},\\
\chi_{\mathrm{CH},ij}&=\delta_{Z_i,6}\delta_{Z_j,1}+\delta_{Z_i,1}\delta_{Z_j,6}.
\end{align}

Current parameters (updated model):
\begin{equation}
K=6.0\times10^{-3},\qquad
\Delta r=0.17~\text{\AA},\qquad
\sigma=0.057~\text{\AA}.
\end{equation}

\paragraph{Note on implementation.}
The practical coded form is:
\begin{equation}
\Delta r_{ij}
=\left[-K\,r^{\mathrm{cov}}_{ij}(1+0.25\,\chi_{\mathrm{CH},ij})N_{ij}f_{\mathrm{coord}}(r)\right]
\left(1-\chi_{\mathrm{CS},ij}e^{-((r-r^{\mathrm{cov}}_{ij}+\Delta r)/\sigma)^2}\right),
\end{equation}
with $\chi_{\mathrm{CH},ij}$ as defined above,
$\chi_{\mathrm{CS},ij}=\chi_{\mathrm{CS}}(i)\chi_{\mathrm{CS}}(j)$,
$N_{ij}=\max(n_i,n_j)^2-1$.

\section{Legacy BDPCS3 (for reference)}
\begin{align}
\Delta r^{\mathrm{legacy}}_{ij} &= \Delta b^{\mathrm{CV}}_{ij}+\Delta b^{\mathrm{V}}_{ij},\\
\Delta b^{\mathrm{CV}}_{ij}
&= -0.0011\left(1+1.1\,\delta_{\mathrm{CH}}\right)\sqrt{n_i n_j-1}\;r^{\mathrm{cov}}_{ij},\\
\Delta b^{\mathrm{V}}_{ij}
&= \Delta b^{\mathrm{CV}}_{ij}\left(\sqrt{|\mathrm{BO}_{ij}-2|}-1\right)
\quad\text{(selected pairs)},
\end{align}
with weak-bond damping:
\begin{equation}
\mathrm{if}\ \mathrm{BO}_{ij}<0.3,\quad \Delta r^{\mathrm{legacy}}_{ij}=0.
\end{equation}
In the unified implementation, $\mathrm{BO}_{ij}$ can also be provided by the topology layer.

\section{Workflow Summary}
\begin{enumerate}
\item Build continuous topology (CNA, effective radii, BO).
\item Build atomic synthons from the same CNA/BO definitions.
\item Use synthons as features for atom typing / chemical-space exploration.
\item Use synthon-consistent BO in BDPCS3 to correct DPCS3 bond lengths.
\item Keep angles/dihedrals unchanged unless additional correction models are introduced.
\end{enumerate}

\section{Outlook}
The same synthon feature space can drive corrections of force-field parameters
(e.g., equilibrium values, force constants, cross terms), enabling transferable,
chemically smooth parameterization across broad chemical spaces.

\end{document}
